% STATUS: written 2026-11-05, recalled results as is (one environment each, source kind and whole statement)
% Sources read in full: Woronowicz 1989 (pp. 125--170, main statements of Sections 1--2).
\section{Quantum groups and the differential calculus of Woronowicz}\label{sec:group-hopf:quantum}

the algebra of Sweedler of Section~\ref{sec:group-hopf:hopf} is the simplest quantum
semigroup. A quantum group in the sense of Woronowicz is a compact matrix
pseudogroup: an algebra $A$ of smooth functions on the group with a comultiplication
\cite[p.~126]{Woronowicz1989}. On a Lie group the exterior derivative satisfies
the Leibniz rule, and the left and right shifts of the group act on the
differential forms \cite[pp.~126--127]{Woronowicz1989}. Woronowicz
takes these two properties as axioms and lets the derivative range over a bimodule:
for a given quantum group there exist many non-equivalent differential calculi
\cite[p.~126]{Woronowicz1989}.

\begin{definition}[{\cite[Definition~1.1, p.~128]{Woronowicz1989}}]\label{def:group-hopf:fodc}
Let $A$ be an algebra with unity, $\Gamma$ be a bimodule over $A$ and
$d\colon A\to\Gamma$ be a linear map. We say that $(\Gamma,d)$ is a
\emph{first order differential calculus} over $A$ if
\begin{enumerate}
\item for any $a,b\in A$, $d(ab)=(da)b+a\,db$;
\item any element $\rho\in\Gamma$ is of the form $\rho=\sum_{k=1}^{K}a_k\,db_k$,
where $a_k,b_k\in A$ and $K$ is a natural number.
\end{enumerate}
\end{definition}

So a first order differential calculus is a derivation into a bimodule, with the
bimodule generated by the image. Let $A^2$ be the kernel of
the multiplication map $m\colon A\otimes A\to A$, made an $A$-bimodule through the
first factor, and set $Db=1\otimes b-b\otimes1$. Then $(A^2,D)$ is a first order
differential calculus over $A$, and it is universal
\cite[p.~128]{Woronowicz1989}.

\begin{proposition}[{\cite[Proposition~1.1, pp.~128--129]{Woronowicz1989}}]\label{prop:group-hopf:universal}
Let $\mathcal{N}$ be a sub-bimodule of $A^2$, $\Gamma=A^2/\mathcal{N}$, $\pi$ be the
canonical epimorphism $A^2\to\Gamma$ and $d=\pi\circ D$. Then $(\Gamma,d)$ is a
first order differential calculus over $A$. Any first order differential calculus
over $A$ can be obtained in this way.
\end{proposition}

So the calculi over a fixed algebra $A$ are the quotients of the universal calculus
$(A^2,D)$. On a quantum group the calculus is asked to be left-covariant: the left
shifts of the group lift to $\Gamma$ \cite[p.~132]{Woronowicz1989}. Woronowicz
classifies such calculi by right ideals of $A$ contained in the kernel of the counit
$\varepsilon$ \cite[p.~132]{Woronowicz1989}.

\begin{theorem}[{\cite[Theorem~1.5, p.~132]{Woronowicz1989}}]\label{thm:group-hopf:classification}
Let $\mathcal{N}$ be a right ideal of $A$ contained in $\ker\varepsilon$ and
$\mathcal{N}'=r^{-1}(A\otimes\mathcal{N})$. Then $\mathcal{N}'$ is a sub-bimodule of
$A^2$. Moreover, let $\Gamma=A^2/\mathcal{N}'$, $\pi$ be the canonical epimorphism
$A^2\to\Gamma$ and $d=\pi\circ D$. Then the first order differential calculus
$(\Gamma,d)$ is left-covariant. Any left-covariant first order differential calculus
on $G$ can be obtained in this way.
\end{theorem}

For a Lie group the choice $\mathcal{N}=(\ker\varepsilon)^2$ recovers the usual
differential calculus \cite[p.~132]{Woronowicz1989}. For covariance under both
shifts the ideal $\mathcal{N}$ has to be ad-invariant
\cite[p.~134]{Woronowicz1989}. A linear subset $T\subset A$ is ad-invariant when
$\operatorname{ad}(T)\subset T\otimes A$, where $\operatorname{ad}$ is the adjoint
coaction \cite[p.~133]{Woronowicz1989}.

\begin{theorem}[{\cite[Theorem~1.8, p.~134]{Woronowicz1989}}]\label{thm:group-hopf:bicovariant}
Let $\mathcal{N}$ be a right ideal of $A$ contained in $\ker\varepsilon$ and
$(\Gamma,d)$ be the left-covariant first order differential calculus described in
Theorem~\ref{thm:group-hopf:classification}. Then $(\Gamma,d)$ is bicovariant if and
only if $\mathcal{N}$ is ad-invariant.
\end{theorem}

The next example is a calculus that is not bicovariant.

\begin{example}[{\cite[p.~132, Example~2; p.~134, remark]{Woronowicz1989}}]\label{ex:group-hopf:3d}
The 3D-calculus on the twisted $SU(2)$ group $SU_\mu(2)$ is left-covariant; its
right ideal $\mathcal{N}$ is generated by six elements, and
$\dim(\ker\varepsilon/\mathcal{N})=3$. This calculus is not bicovariant, while the
standard differential calculus on a Lie group is bicovariant.
\end{example}

A bicovariant calculus then lifts to higher order forms, so the exterior derivative
is well defined on the whole external algebra \cite[p.~125]{Woronowicz1989}. The
classifications of this section all pass through quotient bimodules of $A\otimes A$:
Proposition~\ref{prop:group-hopf:universal} for a bare algebra, and
Theorem~\ref{thm:group-hopf:classification} and Theorem~\ref{thm:group-hopf:bicovariant}
for a quantum group. Section~\ref{sec:group-hopf:crossed-products} returns to the
group itself: a skew derivation of a group ring is a crossed homomorphism on $G$,
and the theorem of Losert makes each of them principal \cite[p.~222]{Losert2008}.

Over a field with a primitive $n$-th root of unity the skew derivations of a finite-dimensional algebra are studied through the action of the double of the Taft Hopf algebra, which is the Hopf-algebraic form of the skew rule~\cite{WlSkewDerivationsOf2001}.

\"Ozav\c{s}ar introduces a two-parameter quantum $(2+1)$-superspace with a Hopf superalgebra structure, defines derivation operators acting on it, and shows that the algebra of these derivations is again a Hopf superalgebra~\cite{WlATwoParameterQuantum21Superspa2013}.
